\section{问题一：【短标题】}

\subsection{模型建立}

【模型假设、变量定义、目标函数与约束。公式编号由 cls 自动生成，
形如 (5.1)——cls 里 \theequation 就是「节号.序号」。】

\begin{equation}
  \label{eq:p1_obj}
  \min \sum_{i=1}^{n} c_i x_i
\end{equation}

\subsection{求解方法}

【算法步骤、停止条件。算法的「输入/输出」标签已由模板设成中文，直接用即可。】

\begin{algorithm}[H]
  \caption{【算法名】}
  \begin{algorithmic}[1]
    \Require 【输入】
    \Ensure  【输出】
    \State 【步骤一】
    \State 【步骤二】
  \end{algorithmic}
\end{algorithm}

\subsection{求解结果}

【关键数值必须与本次求解的实际结果一致，不得编造、估算或另行取整。】

\begin{table}[H]
  \centering
  \caption{【结果表题】}
  \label{tab:p1_result}
  \begin{tabular}{cccc}
    \toprule
    【列一】 & 【列二】 & 【列三】 & 【列四】 \\
    \midrule
    【数据】 & 【数据】 & 【数据】 & 【数据】 \\
    \bottomrule
  \end{tabular}
\end{table}

% 数据图同样只写裸文件名：
% \begin{figure}[H]
%   \centering
%   \includegraphics[width=0.8\textwidth]{fig_q1_error.pdf}
%   \caption{【图题】}
%   \label{fig:p1_error}
% \end{figure}
